\documentclass[10pt,a4paper]{amsart}
\usepackage[margin=19mm]{geometry}
\usepackage{amsmath,amssymb,mathtools}
\usepackage[scr=rsfs,cal=euler]{mathalfa}
\usepackage{xfrac}
\usepackage[unicode,hidelinks,bookmarks=false]{hyperref}
\newcommand{\ie}{i.e.\ }
\newcommand{\cf}{cf.\ }
\newcommand{\resp}{respectively\ }
\newcommand{\Subsection}{\subsection}
\providecommand{\nobreakdash}{\nobreak\hbox{-}}
\setlength{\emergencystretch}{2em}
\multlinegap0pt
\usepackage{amssymb}
\usepackage{xcolor}
\usepackage{tikz}
\usepackage[shortlabels]{enumitem}
\newtheorem{defn}{Definition}
\newtheorem{thm}{Theorem}
\title{Off-Diagonal Ramsey Multiplicity}
\author{Elena Moss, Jonathan A. Noel}
\date{Source: arXiv:2306.17388v2 (CC BY 4.0)}
\begin{document}
\maketitle

\noindent\emph{Source and licence.} Off-Diagonal Ramsey Multiplicity. Version arXiv:2306.17388v2 (CC BY 4.0), distributed by arXiv under the Creative Commons Attribution 4.0 International licence, http://creativecommons.org/licenses/by/4.0/, as recorded in the arXiv licence metadata for this exact version and shown on the abstract page https://arxiv.org/abs/2306.17388v2. Full licence notice: \texttt{licenses/arxiv-2306.17388-CC-BY-4.0.txt}.\par\smallskip

\noindent\textit{Editorial scope.} Counted unit: the article's thm th:commonPair (source0.tex lines 2788--2793). The article's complete original proof of it is delivered with it (source0.tex lines 2795--2863, one \texttt{proof} environment of 69 lines). No mathematics is written, repaired or re-derived: the statement and the proof are the article's own lines, in the article's own notation, and no retained line is substituted or re-flowed.  Retained uncounted as the source-local context the proof needs: lines 2717--2730; lines 2732--2735.  Nothing was omitted from any retained span, so the package carries no omission record.  No retained line contains a citation, so the wrapper carries no bibliography and no bibliography entry.  No retained line contains a figure environment, an image-inclusion command or drawing code, so no image is delivered and the package declares no asset dependency.  Editorial formatting only, in addition: an AMS \texttt{amsart} preamble that restates the article's own declarations and macros used by the retained lines, namely the environments \texttt{defn}, \texttt{thm} on separate counters, together with the standard packages those lines need.  The retained environments print in this document as Definition 1, Theorem 1 in order of appearance, whereas the article prints the same items with the numbering of its own full document.  The complete native source of the article as distributed in the arXiv e-print for this exact version is bundled unchanged as \texttt{sources/143618/source0.tex}.
\begin{defn}
\label{defn:certificate}
Let $H_1$ and $H_2$ be graphs and let $\alpha\in\mathbb{R}$. Say that a pair $(W_1,W_2)$ is an \emph{$\alpha$-certificate} for $(H_1,H_2)$ if there exists $\lambda\in[0,2]$ such that the following four bounds hold:
\begin{enumerate}
\stepcounter{equation}
 \item\label{eq:ralpha} $\lambda\cdot t(H_1,W_1) + (2-\lambda)\cdot t(H_2,1-W_1) \leq \alpha$,
\stepcounter{equation}
 \item\label{eq:balpha} $\lambda\cdot t(H_1,W_2) + (2-\lambda)\cdot t(H_2,1-W_2) \leq \alpha$,
\stepcounter{equation}
 \item\label{eq:rbigger}$t(H_1,W_1)\geq t(H_2,1-W_1)$ and
\stepcounter{equation}
 \item\label{eq:bbigger} $t(H_1,W_2)\leq t(H_2,1-W_2)$.
\end{enumerate}
\end{defn}

\begin{thm}
\label{th:dual}
Let $H_1$ and $H_2$ be non-empty graphs. Then $c(H_1,H_2)$ is equal to the minimum over all $\alpha$ for which there exists an $\alpha$-certificate for $(H_1,H_2)$. 
\end{thm}

\begin{thm}
\label{th:commonPair}
Let $H_1$ and $H_2$ be non-empty graphs and let $p\in(0,1)$ such that $p^{e(H_1)}=(1-p)^{e(H_2)}$. Then
\[c(H_1,H_2)\leq 2\cdot p^{e(H_1)}=2\cdot(1-p)^{e(H_2)}\]
with equality if and only if $(H_1,H_2)$ is $(p,1-p)$-common. 
\end{thm}

\begin{proof}
Let $H_1$ and $H_2$ be arbitrary non-empty graphs. For the upper bound, we show that the pair $(W_1,W_2)$ where $W_1=W_2=p$ is a $2p^{e(H_1)}$-certificate for $(H_1,H_2)$ and apply Theorem~\ref{th:dual}. Since $p^{e(H_1)}=(1-p)^{e(H_2)}$, we have
\[t(H_1,W_1)=t(H_1,W_2)=p^{e(H_1)}=(1-p)^{e(H_2)}=t(H_2,1-W_1)=t(H_2,1-W_2)\]
and so \eqref{eq:rbigger} and \eqref{eq:bbigger} both hold. Using the fact that $p^{e(H_1)}=(1-p)^{e(H_2)}$ again, for any $\lambda\in[0,2]$, we have
\[\lambda\cdot t(H_1,W_1) + (2-\lambda)\cdot t(H_2,1-W_1)\]
\[=\lambda\cdot t(H_1,W_2) + (2-\lambda)\cdot t(H_2,1-W_2)= \lambda\cdot p^{e(H_1)}+ (2-\lambda)(1-p)^{e(H_2)} = 2\cdot p^{e(H_1)}\]
and so \eqref{eq:ralpha} and \eqref{eq:balpha} hold for $\alpha=2p^{e(H_1)}$. Thus, $(W_1,W_2)$ is a $2p^{e(H_1)}$-certificate for $(H_1, H_2)$ and we are done with the upper bound. 

What remains is to characterize the case of equality. Define
\begin{equation}\label{eq:commonlambda}\lambda_0 :=2\left( \frac{e(H_2)(1-p)^{e(H_2)-1}}{e(H_1)p^{e(H_1)-1} + e(H_2)(1-p)^{e(H_2)-1}}\right)\end{equation}
and
\begin{equation}\label{eq:commonsigma}\sigma_0 := \frac{1}{2}\left(\frac{e(H_1)p^{e(H_1)-1} + e(H_2)(1-p)^{e(H_2)-1}}{\left(e(H_1)p^{e(H_1)-1}\right)\cdot\left(e(H_2)(1-p)^{e(H_2)-1}\right)}\right).\end{equation}
Note that
\begin{equation}\label{eq:lambdasigma1}\sigma_0\cdot\lambda_0 = \frac{1}{e(H_1)p^{e(H_1)-1}}\end{equation}
and
\begin{equation}\label{eq:lambdasigma2}\sigma_0\cdot(2-\lambda_0) = \frac{1}{e(H_2)(1-p)^{e(H_2)-1}}.\end{equation}
Therefore, $(H_1,H_2)$ is $(p,1-p)$-common if and only if 
\[\sigma_0\cdot\lambda_0\cdot t(H_1,W) + \sigma_0\cdot(2-\lambda_0)\cdot t(H_2,1-W)\geq \frac{p}{e(H_1)}+\frac{1-p}{e(H_2)}\]
or, equivalently, 
\[\lambda_0\cdot t(H_1,W) + (2-\lambda_0)\cdot t(H_2,1-W)\geq \frac{1}{\sigma_0}\left(\frac{p}{e(H_1)}+\frac{1-p}{e(H_2)}\right)\]
for every graphon $W$. Now, observe that 
\[\frac{1}{\sigma_0}\left(\frac{p}{e(H_1)}+\frac{1-p}{e(H_2)}\right) = \frac{2e(H_1)p^{e(H_1)-1}e(H_2)(1-p)^{e(H_2)-1}}{e(H_1)p^{e(H_1)-1} + e(H_2)(1-p)^{e(H_2)-1}}\left(\frac{p}{e(H_1)}+\frac{1-p}{e(H_2)}\right)\]
\[ = \frac{2e(H_1)p^{e(H_1)-1}e(H_2)(1-p)^{e(H_2)-1}}{e(H_1)p^{e(H_1)-1} + e(H_2)(1-p)^{e(H_2)-1}}\left(\frac{e(H_2)p+e(H_1)(1-p)}{e(H_1)e(H_2)}\right)\]
\[ = \frac{2e(H_2)p^{e(H_1)}(1-p)^{e(H_2)-1}+2e(H_1)p^{e(H_1)-1}(1-p)^{e(H_2)}}{e(H_1)p^{e(H_1)-1} + e(H_2)(1-p)^{e(H_2)-1}}\]
which, since $(1-p)^{e(H_2)}=p^{e(H_1)}$, is equal to $2\cdot p^{e(H_1)}$. So, $(H_1,H_2)$ is $(p,1-p)$-common if and only if 
\[\lambda_0\cdot t(H_1,W) + (2-\lambda_0)\cdot t(H_2,1-W)\geq2\cdot p^{e(H_1)}\]
for every graphon $W$. In particular, if $(H_1,H_2)$ is $(p,1-p)$-common, then
\[c(H_1,H_2)\geq c_{\lambda_0}(H_1,H_2)\geq 2\cdot p^{e(H_1)}\]
which proves one direction of the characterization. 

Now, suppose that $(H_1,H_2)$ is not $(p,1-p)$-common. Our goal is to find an $\alpha$-certificate for $(H_1,H_2)$ with $\alpha<2\cdot p^{e(H_1)}$. Using the characterization above, let $W$ be a graphon such that
\[
\lambda_0\cdot t(H_1,W)+(2-\lambda_0)\cdot t(H_2,1-W)<2\cdot p^{e(H_1)}.
\]
For $\gamma\in[0,2]$, define
\[
g_\gamma(x):=\gamma\cdot (p+x)^{e(H_1)}+(2-\gamma)\cdot(1-p-x)^{e(H_2)}.
\]
By the definition of $\lambda_0$, we have $g'_\gamma(0)>0$ when $\gamma>\lambda_0$ and $g'_\gamma(0)<0$ when $\gamma<\lambda_0$.

First suppose that $t(H_1,W)\geq t(H_2,1-W)$. Choose $\gamma>\lambda_0$ sufficiently close to $\lambda_0$ so that
\[
\gamma\cdot t(H_1,W)+(2-\gamma)\cdot t(H_2,1-W)<2\cdot p^{e(H_1)}.
\]
Set $W_1:=W$ and
\[
\alpha_1:=\gamma\cdot t(H_1,W_1)+(2-\gamma)\cdot t(H_2,1-W_1).
\]
Since $g'_\gamma(0)>0$, we may choose $\varepsilon>0$ sufficiently small so that $p-\varepsilon>0$ and $
g_\gamma(-\varepsilon)<g_\gamma(0)=2p^{e(H_1)}$. Set $W_2:=p-\varepsilon$ and $\alpha_2:=g_\gamma(-\varepsilon)$. Then
\[
t(H_1,W_2)=(p-\varepsilon)^{e(H_1)}<(1-p+\varepsilon)^{e(H_2)}=t(H_2,1-W_2).
\]
Thus $(W_1,W_2)$ is an $\alpha$-certificate with $\alpha:=\max\{\alpha_1,\alpha_2\}<2p^{e(H_1)}$. 

Now suppose instead that $t(H_1,W)<t(H_2,1-W)$. Choose $\gamma<\lambda_0$ sufficiently close to $\lambda_0$ so that
\[
\gamma\cdot t(H_1,W)+(2-\gamma)\cdot t(H_2,1-W)<2\cdot p^{e(H_1)}.
\]
Set $W_2:=W$ and
\[
\alpha_2:=\gamma\cdot t(H_1,W_2)+(2-\gamma)\cdot t(H_2,1-W_2).
\]
Since $g'_\gamma(0)<0$, we may choose $\varepsilon>0$ sufficiently small so that $p+\varepsilon<1$ and $g_\gamma(\varepsilon)<g_\gamma(0)=2p^{e(H_1)}$. Set $W_1:=p+\varepsilon$ and $\alpha_1:=g_\gamma(\varepsilon)$. Then
\[
t(H_1,W_1)=(p+\varepsilon)^{e(H_1)}>(1-p-\varepsilon)^{e(H_2)}=t(H_2,1-W_1).
\]
Thus $(W_1,W_2)$ is an $\alpha$-certificate with $\alpha:=\max\{\alpha_1,\alpha_2\}<2p^{e(H_1)}$. In both cases, Theorem~\ref{th:dual} gives $c(H_1,H_2)<2p^{e(H_1)}$, as desired.
\end{proof}

\end{document}
